\documentclass{article}

\usepackage[utf8]{inputenc}
\usepackage[T1]{fontenc}
\usepackage[a4paper,margin=2.5cm]{geometry}
\usepackage{xcolor}
\usepackage{amsmath}
\usepackage{amssymb}
\usepackage{listings}

\lstset{
  basicstyle=\ttfamily\small,
  breaklines=true,
  frame=single,
  numbers=left,
  numberstyle=\tiny\color{gray},
  keywordstyle=\bfseries\color{blue!60!black},
  commentstyle=\itshape\color{green!40!black},
  stringstyle=\color{red!50!black},
  showstringspaces=false,
  tabsize=2
}

\title{Sistemi di Tipi, Sicurezza della Memoria e Garbage Collection}
\author{}
\date{}

\begin{document}

\maketitle

\section{Relazioni di Preordine ed Equivalenza}

Una relazione binaria $*$ su un insieme $D$ è un sottoinsieme del prodotto cartesiano:
\[
* \subseteq D \times D.
\]
Scriviamo $c * d$ per indicare $(c,d) \in *$.

Una relazione è:
\begin{itemize}
  \item \textbf{riflessiva} se $\forall d \in D,\ d*d$;
  \item \textbf{simmetrica} se $\forall c,d \in D,\ c*d \Rightarrow d*c$;
  \item \textbf{transitiva} se $\forall c,d,e \in D,\ (c*d \land d*e) \Rightarrow c*e$;
  \item \textbf{antisimmetrica} se $\forall c,d \in D,\ (c*d \land d*c) \Rightarrow c=d$.
\end{itemize}

Una relazione riflessiva e transitiva è un \textbf{preordine}. Un preordine simmetrico è una \textbf{equivalenza}; un preordine antisimmetrico è un \textbf{ordine parziale}.

Nel contesto dei \textbf{sistemi di tipi}, le definizioni di tipo possono essere interpretate in due modi:
\begin{itemize}
  \item come definizioni \textbf{opache}, che conducono all'\textbf{equivalenza per nome};
  \item come definizioni \textbf{trasparenti}, che conducono all'\textbf{equivalenza strutturale}.
\end{itemize}

\section{Equivalenza tra Tipi}

\subsection{Equivalenza per Nome}

\textbf{Definizione.} Due tipi sono equivalenti per nome se e solo se hanno lo stesso nome. In questa interpretazione, ogni nuova definizione introduce un nuovo tipo distinto.

\begin{lstlisting}
type T1 = 1..10;
type T2 = 1..10;
type T3 = int;
type T4 = int;
\end{lstlisting}

Con \textbf{equivalenza per nome forte}, \texttt{T1}, \texttt{T2}, \texttt{T3} e \texttt{T4} sono quattro tipi distinti. Con \textbf{equivalenza per nome debole}, una semplice rinomina non introduce un nuovo tipo: \texttt{T3} e \texttt{T4} possono essere considerati alias dello stesso tipo, mentre \texttt{T1} e \texttt{T2} restano distinti.

Un problema tipico emerge con i tipi anonimi. In Pascal, dichiarazioni strutturalmente identiche possono produrre tipi diversi:

\begin{lstlisting}
var V: array [1..10] of integer;
    W: array [1..10] of integer;

procedure f(Z: array [1..10] of integer);
\end{lstlisting}

Con equivalenza per nome, \texttt{V}, \texttt{W} e il parametro formale \texttt{Z} appartengono a tre tipi anonimi distinti.

\subsection{Equivalenza Strutturale}

\textbf{Definizione.} L'\textbf{equivalenza strutturale} è la più piccola relazione di equivalenza sui tipi tale che:
\begin{enumerate}
  \item il nome di un tipo è equivalente a sé stesso;
  \item se un tipo $T$ è introdotto da una definizione \texttt{type T = expression}, allora $T$ è equivalente a \texttt{expression};
  \item se due tipi sono costruiti applicando lo stesso costruttore di tipo a tipi equivalenti, allora i due tipi risultanti sono equivalenti.
\end{enumerate}

\begin{lstlisting}
type T1 = int;
type T2 = char;
type T3 = struct{
    T1 a;
    T2 b;
}
type T4 = struct{
    int a;
    char b;
}
\end{lstlisting}

\texttt{T3} e \texttt{T4} sono strutturalmente equivalenti, poiché dopo l'espansione dei nomi hanno la stessa struttura.

L'equivalenza strutturale solleva però questioni sottili sui nomi dei campi, sull'ordine dei campi e sui tipi ricorsivi:

\begin{lstlisting}
type S = struct{
    int a;
    int b;
}
type T = struct{
    int n;
    int m;
}
type U = struct{
    int m;
    int n;
}
\end{lstlisting}

\begin{lstlisting}
type R1 = struct{
    int a;
    R2 p;
}
type R2 = struct{
    int a;
    R1 p;
}
\end{lstlisting}

Nel caso dei tipi ricorsivi, l'equivalenza può richiedere argomenti matematici più sofisticati della semplice espansione sintattica.

\subsection{Scelte nei Linguaggi}

I linguaggi reali combinano spesso le due nozioni:
\begin{itemize}
  \item Pascal usa equivalenza per nome debole;
  \item Java usa equivalenza per nome, ma equivalenza strutturale per gli array;
  \item C usa equivalenza strutturale per array e \texttt{typedef}, ma equivalenza per nome per \texttt{struct} e \texttt{union};
  \item C\texttt{++} usa prevalentemente equivalenza per nome;
  \item ML usa equivalenza strutturale, salvo per tipi definiti con \texttt{datatype}.
\end{itemize}

\section{Compatibilità e Conversione}

\subsection{Compatibilità}

\textbf{Definizione.} Un tipo $T$ è \textbf{compatibile} con un tipo $S$ se un valore di tipo $T$ può essere usato in ogni contesto in cui sarebbe ammissibile un valore di tipo $S$.

La compatibilità è più debole dell'equivalenza. Se due tipi sono equivalenti, allora sono compatibili; il viceversa non vale in generale. La compatibilità è tipicamente un \textbf{preordine}: è riflessiva e transitiva, ma non necessariamente simmetrica.

La compatibilità può essere interpretata con diversi gradi di permissività:
\begin{enumerate}
  \item $T$ e $S$ sono equivalenti;
  \item i valori di $T$ sono un sottoinsieme dei valori di $S$;
  \item tutte le operazioni ammesse su $S$ sono ammesse anche su $T$;
  \item i valori di $T$ corrispondono canonicamente a valori di $S$;
  \item i valori di $T$ possono essere trasformati convenzionalmente in valori di $S$.
\end{enumerate}

Esempio di compatibilità strutturale per record:

\begin{lstlisting}
type S = struct{
    int a;
}
type T = struct{
    int a;
    char b;
}
\end{lstlisting}

Ogni operazione su \texttt{S} che seleziona il campo \texttt{a} ha senso anche su valori di tipo \texttt{T}. Questa intuizione anticipa la nozione di \textbf{sottotipo}.

\subsection{Coercion e Cast}

Una \textbf{conversione implicita}, detta anche \textbf{coercion}, è inserita dalla macchina astratta senza annotazioni esplicite nel programma. Una \textbf{conversione esplicita}, o \textbf{cast}, è invece indicata direttamente dal programmatore.

Nel caso generale:
\[
t : T,\qquad T \text{ compatibile con } S
\]
allora un contesto che richiede $S$ può accettare $t$, eventualmente dopo una conversione:
\[
\operatorname{coerce}_{T \to S}(t) : S.
\]

I cast possono anche non modificare la rappresentazione in memoria. Questo è pericoloso nei linguaggi non type-safe:

\begin{lstlisting}[language=C]
int a = 233;
float b = *((float*) &a);
\end{lstlisting}

Qui l'indirizzo di \texttt{a} viene reinterpretato come puntatore a \texttt{float}. I bit non cambiano, ma cambia il tipo con cui vengono letti. Questo tipo di conversione è una possibile fonte di errori difficili da individuare.

Nella notazione del capitolo, un cast è scritto come:

\begin{lstlisting}
S s = (S) t;
\end{lstlisting}

I linguaggi moderni tendono a preferire i \textbf{cast espliciti} alle coercion implicite, perché documentano meglio l'intenzione del programmatore e interagiscono in modo più prevedibile con overloading e polimorfismo.

\section{Polimorfismo}

\subsection{Nozione Generale}

Un sistema di tipi è \textbf{monomorfico} se ogni oggetto del linguaggio ha un unico tipo. È invece \textbf{polimorfico} se lo stesso oggetto può avere più di un tipo.

\textbf{Definizione.} Un sistema di tipi in cui lo stesso oggetto può avere più di un tipo è detto \textbf{polimorfico}; per estensione, un oggetto è polimorfico quando il sistema di tipi gli assegna più di un tipo.

Esempi ricorrenti:
\[
+ : int \times int \to int,
\qquad
+ : float \times float \to float.
\]

Il valore \texttt{null} può avere tipo $T^*$ per ogni tipo $T$, e una funzione \texttt{length} su array può avere tipo:
\[
T[] \to int.
\]

\subsection{Motivazione}

Senza polimorfismo, funzioni identiche devono essere duplicate cambiando solo le annotazioni di tipo:

\begin{lstlisting}
void int_sort(int A[])
\end{lstlisting}

\begin{lstlisting}
void char_sort(char C[])
\end{lstlisting}

In un linguaggio polimorfico si può scrivere un'unica funzione generica:

\begin{lstlisting}
void sort(<T> A[])
\end{lstlisting}

Le principali forme di polimorfismo sono:
\begin{itemize}
  \item \textbf{polimorfismo ad hoc}, o \textbf{overloading};
  \item \textbf{polimorfismo universale parametrico};
  \item \textbf{polimorfismo universale per sottotipo}.
\end{itemize}

\subsection{Overloading}

Un nome è \textbf{overloaded} quando denota più oggetti distinti e il contesto deve determinare quale oggetto è effettivamente inteso.

Esempi:
\begin{itemize}
  \item \texttt{+} come somma intera, somma reale o concatenazione;
  \item più funzioni con lo stesso nome ma parametri diversi.
\end{itemize}

L'overloading viene risolto staticamente usando informazioni di tipo. Concettualmente, il compilatore sostituisce ogni nome ambiguo con un nome non ambiguo.

L'interazione tra overloading e coercion può rendere ambigua la semantica:

\begin{lstlisting}
1 + 2
1.0 + 2.0
1 + 2.0
1.0 + 2
\end{lstlisting}

Il linguaggio deve specificare se \texttt{+} ha molteplici significati, se vengono inserite coercion, oppure se entrambe le tecniche cooperano.

\subsection{Polimorfismo Parametrico Universale}

\textbf{Definizione.} Un valore mostra \textbf{polimorfismo parametrico universale} quando possiede un numero infinito di tipi ottenibili istanziando un unico schema di tipo generale.

Una funzione parametricamente polimorfica è implementata da un solo algoritmo uniforme, che non dipende dalla struttura specifica del tipo.

La notazione:
\[
\forall T.\ T[] \to void
\]
esprime una funzione applicabile ad array di qualunque tipo $T$.

Esempio: scambio di due variabili dello stesso tipo generico.

\begin{lstlisting}
void swap(reference <T> x, reference <T> y){
    <T> tmp = x;
    x = y;
    y = tmp;
}
\end{lstlisting}

Esempio di istanziazione automatica:

\begin{lstlisting}
int* k = null;
char v, w;
int i, j;

swap(v, w);
swap(i, j);
\end{lstlisting}

Nel primo caso \texttt{null} è istanziato come \texttt{int*}; poi \texttt{swap} viene istanziata rispettivamente sui caratteri e sugli interi.

\subsection{Polimorfismo Esplicito e Implicito}

Nel \textbf{polimorfismo esplicito}, il programma contiene annotazioni visibili dei parametri di tipo, come \texttt{<T>}. È il modello dei template di C\texttt{++} e dei generics di Java.

Nel \textbf{polimorfismo implicito}, tipico di ML, il programmatore non fornisce annotazioni di tipo; il sistema inferisce il tipo più generale.

Esempio di funzione identità:

\begin{lstlisting}
fun Ide(x){return x;}
\end{lstlisting}

Il tipo inferito è:
\[
\forall T.\ T \to T.
\]

Esempio di composizione:

\begin{lstlisting}
fun Comp(f, g, x){return f(g(x));}
\end{lstlisting}

Il tipo più generale è:
\[
(S \to T) \times (R \to S) \times R \to T.
\]

In forma polimorfica:
\[
\forall R,S,T.\ (S \to T) \times (R \to S) \times R \to T.
\]

\subsection{Esempi in ML}

\begin{lstlisting}
- fun Ide(x) = x;
val Ide = fn : 'a -> 'a

- Ide(4);
val it = 4 : int

- Ide(true);
val it = true : bool
\end{lstlisting}

\begin{lstlisting}
- fun Comp(f,g,x) = f(g(x));
val Comp = fn : ('a -> 'b) * ('c -> 'a) * 'c -> 'b
\end{lstlisting}

\begin{lstlisting}
- fun swap(x,y) = let tmp = !x in
                    (x := !y; y := tmp;)
                  end;
val swap = fn : ('a ref) * ('a ref) -> unit
\end{lstlisting}

Il tipo \texttt{'a ref} indica una variabile o riferimento di tipo arbitrario. Il tipo \texttt{unit} è il tipo singleton usato per espressioni con effetti collaterali.

\subsection{Polimorfismo per Sottotipo}

Il \textbf{polimorfismo per sottotipo} è tipico dei linguaggi orientati agli oggetti. Supponendo una relazione di sottotipo $<:$, la scrittura:
\[
C <: D
\]
si legge ``$C$ è sottotipo di $D$''.

\textbf{Definizione.} Un valore mostra \textbf{polimorfismo per sottotipo} quando esistono infiniti tipi ottenibili istanziando uno schema generale con i sottotipi di un tipo assegnato.

Una funzione di tipo:
\[
\forall T <: D.\ T \to void
\]
può essere applicata a valori di qualunque sottotipo di $D$. Il polimorfismo è dunque limitato dal vincolo di sottotipo.

\subsection{Implementazione del Polimorfismo}

Un primo approccio, tipico di C\texttt{++}, risolve il polimorfismo staticamente, generando una copia del codice per ogni istanziazione.

\begin{lstlisting}[language=C++]
template<typename T>
void swap(T& x, T& y){
    T tmp = x;
    x = y;
    y = tmp;
}
\end{lstlisting}

Questo produce codice efficiente a runtime, ma può aumentare la dimensione dell'eseguibile.

Un secondo approccio mantiene una sola versione del codice polimorfico. I dati vengono rappresentati indirettamente tramite puntatori e descrittori dinamici. Questa soluzione riduce la duplicazione del codice, ma introduce un costo di accesso indiretto.

\section{Type Checking e Type Inference}

\subsection{Type Checking}

Il \textbf{type checker} verifica che il programma rispetti le regole imposte dal sistema di tipi, in particolare le regole di compatibilità.

Nel \textbf{controllo statico}, il type checker è un modulo del compilatore. Nel \textbf{controllo dinamico}, è parte del sistema runtime.

Il type checker visita l'\textbf{albero di sintassi astratta} dal basso verso l'alto:
\begin{itemize}
  \item parte da variabili e costanti, i cui tipi sono noti;
  \item calcola il tipo delle espressioni composte;
  \item verifica che operatori, assegnamenti e chiamate siano coerenti con il sistema di tipi.
\end{itemize}

Esempio:

\begin{lstlisting}
int f(int n){return n + 1;}
\end{lstlisting}

Dal fatto che \texttt{n} è \texttt{int}, che \texttt{1} è \texttt{int} e che \texttt{+} applicato a due interi restituisce un intero, il sistema può verificare che il risultato è \texttt{int}.

\subsection{Type Inference}

La \textbf{type inference} attribuisce un tipo a un'espressione anche in assenza di dichiarazioni esplicite.

\begin{lstlisting}
fun f(n){return n + 1;}
\end{lstlisting}

Poiché \texttt{1} è intero e \texttt{+} richiede due interi, l'algoritmo inferisce:
\[
f : int \to int.
\]

Schema semplificato dell'inferenza:
\begin{enumerate}
  \item assegnare un tipo a ogni nodo dell'albero sintattico; per nomi non ancora noti usare variabili di tipo;
  \item attraversare l'albero generando vincoli di uguaglianza tra tipi;
  \item risolvere i vincoli mediante \textbf{unificazione}.
\end{enumerate}

Se una funzione $f$ ha tipo iniziale $\alpha$ e viene applicata a un argomento $v$ di tipo $\beta$, si genera il vincolo:
\[
\alpha = \beta \to \gamma
\]
dove $\gamma$ è una nuova variabile di tipo per il risultato.

Esempio:

\begin{lstlisting}
fun g(n){return n;}
\end{lstlisting}

Il tipo inferito è:
\[
\forall T.\ T \to T.
\]

La presenza di variabili di tipo non risolte non è un errore: indica una espressione polimorfica.

\section{Safety dei Linguaggi}

La \textbf{safety} basata sui tipi riguarda la garanzia che l'esecuzione di un programma ben tipato non produca errori nascosti dovuti a violazioni di tipo.

I linguaggi possono essere classificati in:
\begin{enumerate}
  \item \textbf{non safe};
  \item \textbf{localmente safe};
  \item \textbf{safe}.
\end{enumerate}

\subsection{Linguaggi Non Safe}

Un linguaggio è \textbf{non safe} quando il sistema di tipi può essere aggirato dal programmatore. Appartengono a questa categoria i linguaggi che permettono:
\begin{itemize}
  \item accesso diretto alla rappresentazione dei dati;
  \item aritmetica sui puntatori;
  \item reinterpretazioni non controllate della memoria.
\end{itemize}

C e C\texttt{++} sono esempi classici di linguaggi non safe.

\subsection{Linguaggi Localmente Safe}

I linguaggi \textbf{localmente safe} hanno un sistema di tipi regolato e controllato, ma includono costrutti limitati che possono introdurre insicurezza. Esempi sono Algol, Pascal e Ada, quando dispongono di controlli runtime adeguati.

Le principali fonti di insicurezza sono:
\begin{itemize}
  \item \textbf{unioni} o varianti non controllate;
  \item \textbf{deallocazione esplicita} della memoria.
\end{itemize}

Esempio in Pascal di variante che permette di manipolare un puntatore come intero:

\begin{lstlisting}
var v: record
    case bool of
        true  : (i: integer);
        false : (p: ^integer)
end
\end{lstlisting}

\subsection{Linguaggi Safe}

Un linguaggio è \textbf{safe} se un teorema garantisce che l'esecuzione di un programma tipato non può generare errori nascosti indotti da violazioni di tipo.

Esempi:
\begin{itemize}
  \item Lisp e Scheme, con controllo dinamico dei tipi;
  \item ML, con controllo statico;
  \item Java, con controllo statico e numerosi controlli runtime.
\end{itemize}

\section{Dangling References}

Una \textbf{dangling reference} è un riferimento a un oggetto la cui vita è terminata. Dereferenziare tale riferimento può produrre errori di sicurezza e corruzione della memoria.

La macchina astratta può includere meccanismi dinamici per prevenire la dereferenziazione di puntatori pendenti.

\subsection{Tombstone}

Con la tecnica delle \textbf{tombstone}, a ogni oggetto allocato viene associata una parola aggiuntiva di memoria, detta tombstone.

Il puntatore non contiene direttamente l'indirizzo dell'oggetto, ma l'indirizzo della tombstone. La tombstone contiene l'indirizzo dell'oggetto:
\[
\text{puntatore} \to \text{tombstone} \to \text{oggetto}.
\]

Quando l'oggetto viene deallocato, la tombstone viene invalidata con un valore speciale. Ogni successiva dereferenziazione passa dalla tombstone e può rilevare l'errore.

Vantaggi:
\begin{itemize}
  \item rileva ogni tentativo di dereferenziare un riferimento pendente;
  \item funziona sia per puntatori allo heap sia, con opportune cautele, per puntatori allo stack.
\end{itemize}

Svantaggi:
\begin{itemize}
  \item richiede un secondo livello di indirezione;
  \item consuma memoria aggiuntiva;
  \item le tombstone invalide possono accumularsi;
  \item può richiedere un garbage collector per riciclare tombstone non più referenziate.
\end{itemize}

\subsection{Locks and Keys}

La tecnica \textbf{locks and keys} associa a ogni oggetto dello heap una \textbf{lock}. Ogni puntatore contiene:
\[
(\text{indirizzo},\ \text{key}).
\]

La \textbf{key} viene inizializzata con il valore della lock dell'oggetto. A ogni dereferenziazione, la macchina astratta controlla:
\[
\text{key} = \text{lock}.
\]

Se il controllo fallisce, viene segnalato un errore. Quando l'oggetto viene deallocato, la lock è invalidata.

Rispetto alle tombstone:
\begin{itemize}
  \item non accumula strutture invalide permanenti;
  \item richiede più spazio per i puntatori, perché ogni puntatore contiene anche una key;
  \item introduce un costo su assegnamento e dereferenziazione dei puntatori;
  \item offre una protezione probabilistica nel caso di riuso della memoria.
\end{itemize}

\section{Garbage Collection}

Nei linguaggi senza deallocazione esplicita, la macchina astratta deve recuperare automaticamente la memoria heap non più utilizzata. Tale compito è svolto dal \textbf{garbage collector}.

Dal punto di vista logico, un garbage collector ha due fasi:
\begin{enumerate}
  \item \textbf{garbage detection}: distinguere gli oggetti vivi da quelli non più utilizzabili;
  \item \textbf{reclamation}: recuperare la memoria degli oggetti non più in uso.
\end{enumerate}

In pratica, le due fasi possono essere intrecciate. Inoltre, la nozione di oggetto ``non più in uso'' è spesso una approssimazione conservativa.

\subsection{Root Set e Descrittori}

Le tecniche di garbage collection devono riconoscere quali campi di un oggetto sono puntatori. Se i tipi sono statici, il compilatore può generare un \textbf{descrittore di tipo} contenente gli offset dei campi puntatore.

Il \textbf{root set} è l'insieme dei riferimenti attivi da cui parte la visita degli oggetti raggiungibili, tipicamente:
\begin{itemize}
  \item puntatori presenti nello stack;
  \item registri;
  \item variabili globali.
\end{itemize}

\subsection{Reference Counting}

Nel \textbf{reference counting}, a ogni oggetto heap è associato un contatore che registra il numero di puntatori attivi verso quell'oggetto.

Alla creazione:
\[
rc(o) = 1.
\]

Per un assegnamento:

\begin{lstlisting}
p = q;
\end{lstlisting}

la macchina astratta:
\begin{itemize}
  \item incrementa il contatore dell'oggetto puntato da \texttt{q};
  \item decrementa il contatore dell'oggetto precedentemente puntato da \texttt{p}.
\end{itemize}

Quando:
\[
rc(o)=0,
\]
l'oggetto $o$ può essere deallocato. Se contiene puntatori interni, i contatori degli oggetti puntati vengono decrementati ricorsivamente.

Vantaggi:
\begin{itemize}
  \item è semplice;
  \item è incrementale;
  \item può essere adatto a sistemi real-time.
\end{itemize}

Limiti:
\begin{itemize}
  \item non recupera strutture circolari irraggiungibili;
  \item introduce costo su assegnamenti e uscita dagli ambienti locali;
  \item il costo dipende dal lavoro del programma, non solo dalla dimensione dello heap.
\end{itemize}

\subsection{Mark and Sweep}

Il metodo \textbf{mark and sweep} implementa esplicitamente le due fasi:

\begin{enumerate}
  \item \textbf{Mark}: tutti gli oggetti sono inizialmente marcati come non in uso; partendo dal root set, si visitano ricorsivamente gli oggetti raggiungibili e li si marca come in uso.
  \item \textbf{Sweep}: lo heap viene scandito; gli oggetti marcati come in uso restano allocati, gli altri tornano nella free list.
\end{enumerate}

Limiti principali:
\begin{itemize}
  \item non è incrementale;
  \item può causare pause percepibili;
  \item richiede tempo proporzionale alla dimensione totale dello heap;
  \item non elimina la frammentazione esterna;
  \item può peggiorare la località di riferimento.
\end{itemize}

\subsection{Pointer Reversal}

La fase di marking richiede una visita ricorsiva del grafo degli oggetti. Se la memoria è quasi esaurita, usare uno stack ausiliario può essere problematico.

La tecnica del \textbf{pointer reversal} usa temporaneamente i puntatori della struttura visitata per rappresentare lo stack di visita. Al termine, i puntatori vengono ripristinati:
\[
\text{struttura finale} = \text{struttura iniziale} \quad \text{a meno dei marchi di visita}.
\]

Questa tecnica consente di visitare strutture collegate usando solo un numero costante di puntatori ausiliari, come \texttt{p\_prec} e \texttt{p\_curr}.

\subsection{Mark and Compact}

Il metodo \textbf{mark and compact} sostituisce la fase di sweep con una fase di compattazione. Gli oggetti vivi sono spostati in modo contiguo, lasciando un unico blocco di memoria libera.

La compattazione richiede tipicamente:
\begin{enumerate}
  \item calcolo della nuova posizione degli oggetti vivi;
  \item aggiornamento dei puntatori interni;
  \item spostamento effettivo degli oggetti.
\end{enumerate}

Vantaggi:
\begin{itemize}
  \item elimina la frammentazione esterna;
  \item migliora la località;
  \item rende la memoria libera gestibile come un unico blocco.
\end{itemize}

Svantaggio principale: è più costoso di mark and sweep quando molti oggetti devono essere spostati.

\subsection{Stop and Copy}

Nel garbage collector \textbf{stop and copy}, lo heap è diviso in due semi-spazi:
\begin{itemize}
  \item \textbf{fromspace}, usato durante l'esecuzione normale;
  \item \textbf{tospace}, usato durante la raccolta.
\end{itemize}

Quando il fromspace si esaurisce, il collector parte dal root set, copia nel tospace tutti gli oggetti vivi e li compatta. Al termine, i ruoli dei due semi-spazi vengono scambiati.

Il costo della raccolta è proporzionale al numero di oggetti vivi:
\[
T_{\mathrm{GC}} = O(\#\text{oggetti vivi}).
\]

L'allocazione normale è molto efficiente, poiché avviene per semplice avanzamento di un puntatore.

\subsection{Algoritmo di Cheney}

L'\textbf{algoritmo di Cheney} realizza efficientemente stop and copy trattando gli oggetti copiati nel tospace come una coda:
\begin{enumerate}
  \item copia nel tospace gli oggetti immediatamente raggiungibili dal root set;
  \item considera il primo oggetto della coda;
  \item copia in fondo alla coda gli oggetti puntati dai suoi campi;
  \item aggiorna i puntatori verso le nuove copie;
  \item continua finché la coda è vuota.
\end{enumerate}

Al termine, il tospace contiene una copia compatta di tutti e soli gli oggetti vivi.

\section{Quadro Riassuntivo}

\begin{center}
\begin{tabular}{|p{3.2cm}|p{5.2cm}|p{5.2cm}|}
\hline
\textbf{Tecnica} & \textbf{Idea centrale} & \textbf{Limite principale} \\
\hline
Equivalenza per nome & Due tipi coincidono solo se hanno lo stesso nome & Può essere troppo restrittiva con tipi anonimi \\
\hline
Equivalenza strutturale & Due tipi coincidono se hanno la stessa struttura & Ambigua con campi nominati, ordine e tipi ricorsivi \\
\hline
Coercion & Conversione implicita inserita dalla macchina astratta & Può nascondere trasformazioni costose o pericolose \\
\hline
Cast & Conversione esplicita scritta dal programmatore & Può violare la safety se reinterpreta la memoria \\
\hline
Polimorfismo parametrico & Un algoritmo uniforme vale per molti tipi & Implementazione potenzialmente costosa o duplicata \\
\hline
Polimorfismo per sottotipo & Istanziazione limitata ai sottotipi di un tipo dato & Dipende da una relazione di sottotipo ben definita \\
\hline
Tombstone & Puntatori indiretti tramite celle invalidabili & Alto costo di spazio e tempo \\
\hline
Locks and keys & Puntatore come coppia indirizzo-key & Costo su puntatori e controllo probabilistico \\
\hline
Reference counting & Ogni oggetto conta i puntatori entranti & Non raccoglie cicli \\
\hline
Mark and sweep & Marca i raggiungibili e libera il resto & Frammentazione e pause non incrementali \\
\hline
Mark and compact & Sposta i vivi in memoria contigua & Più passate e costo di spostamento \\
\hline
Stop and copy & Copia solo gli oggetti vivi in un semi-spazio & Richiede memoria divisa in due semi-spazi \\
\hline
\end{tabular}
\end{center}

\section{Sintesi Finale}

Il capitolo mostra che la \textbf{sicurezza della memoria} e la \textbf{sicurezza dei tipi} sono strettamente connesse. Le nozioni di equivalenza, compatibilità, conversione, polimorfismo e inferenza determinano quali programmi sono considerati corretti dal linguaggio; i meccanismi runtime, come tombstone, locks and keys e garbage collection, proteggono invece l'esecuzione concreta da errori legati a riferimenti pendenti e gestione dello heap.

Un linguaggio è tanto più sicuro quanto più impedisce al programmatore di violare l'astrazione dei tipi e quanto più la macchina astratta può garantire che ogni accesso alla memoria sia coerente con le regole del sistema di tipi.

\end{document}
